\documentclass[12 pt]{extarticle}

	\usepackage[frenchb]{babel}
	\usepackage[utf8]{inputenc}  
	\usepackage[T1]{fontenc}
	\usepackage{amssymb}
	\usepackage[mathscr]{euscript}
	\usepackage{stmaryrd}
	\usepackage{amsmath}
	\usepackage{tikz}
	\usepackage[all,cmtip]{xy}
	\usepackage{amsthm}
	\usepackage{varioref}
	\usepackage{geometry}
	\geometry{a4paper}
	\usepackage{lmodern}
	\usepackage{multicol}
	\usepackage{hyperref}
	\usepackage{array}
	 \usepackage{fancyhdr}
\renewcommand{\theenumi}{\alph{enumi})}
	\pagestyle{fancy}
	\theoremstyle{plain}
	\fancyfoot[C]{} 
	\fancyhead[L]{Contrôle}
	\fancyhead[R]{29 septembre 2025}\geometry{
 a4paper,
 total={170mm,257mm},
 left=10mm,right=10mm,
 top=5mm, bottom = 10 mm
 }
	
	
	\title{Interrogation Chapitre 2}
	\date{}
	\begin{document}

%\begin{center}{\Large Interrogation}\\ 
% \end{center}
 
 \subsection*{Exercice 1 (6 points)}
\begin{enumerate}
\item Écrire $1954,425$ comme une fraction décimale.
\item Écrire $953,238$ sous forme mixte. 
\item Écrire $\frac{43}{1~000}$ sous forme décimale.
\item Écrire $\frac{54~133}{100}$ sous forme mixte. 
\item Donner l'écriture décimale de $2\times 100 + 4 + \frac3{10} + \frac7{1~000}$
\item Donner l'écriture décimale de $\frac{12}{100}+ \frac{32}{1~000}$.
\end{enumerate}
 \subsection*{Exercice 2 (6 points)}
 % 1+ 1,5 + ,5 + 1 + 2= 6
 \textbf{Tous les tracés de cet exercice sont sur la même figure.}
 \begin{enumerate}
 \item Construire un segment d'extrémités $A$ et $B$. 
 \item Construire la perpendiculaire à $(AB)$ passant par $B$.  
 \item Placer un point $C$ sur cette droite.
 \item Tracer la demi-droite d'origine $C$ passant par $A$.
 \item Donner les noms en notation mathématique du segment, de la droite, et de la demi-droite tracés aux questions $a).$, $c).$ et $d).$.
	
 
 
 \end{enumerate}
 \subsection*{Exercice 3 (8 points)}
 \includegraphics[scale=.5]{Exercice-droites}
 \begin{enumerate}
 \item Nommer deux droites parallèles de cette figure. 
 \item Nommer deux droites sécantes de cette figure. 
 \item Donner le point d'intersection des deux droites précédentes. 
 \item Nommer trois droites concourantes. 
 \item \textbf{Recopier} et compléter avec les symboles mathématiques appropriés:  
 \[ F \ldots [CD] \phantom{mioaouoreuououz} G \ldots [FD) \]
 \[ E \ldots (FA] \phantom{mioaouoreuououz} (DF) \ldots (BC) \]
 \end{enumerate}
  
\newpage
 
 
 \subsection*{Exercice 1 (6 points)}
\begin{enumerate}
\item Écrire $185,443$ comme une fraction décimale.
\item Écrire $713,238$ sous forme mixte. 
\item Écrire $\frac{34}{1~000}$ sous forme décimale.
\item Écrire $\frac{9~133}{100}$ sous forme mixte. 
\item Donner l'écriture décimale de $3\times 10 + 4 + \frac3{100} + \frac7{1~000}$
\item Donner l'écriture décimale de $\frac{53}{100}+ \frac{32}{1~000}$.
\end{enumerate}
 \subsection*{Exercice 2 (6 points)}
 % 1+ 1,5 + ,5 + 1 + 2= 6
 \textbf{Tous les tracés de cet exercice sont sur la même figure.}
 \begin{enumerate}
 \item Construire un segment d'extrémités $A$ et $B$. 
 \item Construire la perpendiculaire à $(AB)$ passant par $B$.  
 \item Placer un point $C$ sur cette droite.
 \item Tracer la demi-droite d'origine $C$ passant par $A$.
 \item Donner les noms en notation mathématique du segment, de la droite, et de la demi-droite tracés aux questions $a).$, $c).$ et $d).$.
	
 
 
 \end{enumerate}
 \subsection*{Exercice 3 (8 points)}
 \includegraphics[scale=.5]{Exercice-droites}
 \begin{enumerate}
 \item Nommer deux droites sécantes de cette figure. 
 \item Donner le point d'intersection des deux droites précédentes. 
 \item Nommer deux droites parallèles de cette figure. 
 \item Nommer trois droites concourantes. 
 \item \textbf{Recopier} et compléter avec les symboles mathématiques appropriés:  
 \[ F \ldots [EB] \phantom{mioaouoreuououz} G \ldots [AH) \]
 \[ E \ldots (FA] \phantom{mioaouoreuououz} (DF) \ldots (BC) \]
 \end{enumerate}

 
 	\end{document}
